\textbf{Proof.}
Let \(q_a=F_a^{-1}\) and \(\widehat q_a=\widehat F_a^{-1}\).  For independent centered variables bounded by one, fourth-moment expansion gives
\[
\mathbb E\left(\sum_{j=1}^mX_j\right)^4\le3m^2.
\tag{1}
\]
Markov and Borel--Cantelli along \(m=k^2\), followed by bounded interpolation, prove the strong law for bounded pilot functions.  Apply it to \(y\), \(y^2\), and \(\mathbf1\{y\le t\}\) for rational \(t\).  Monotonicity and rational squeezing give weak empirical convergence.  On \([0,1]\), the area identity
\[
\int_0^1|R^{-1}(u)-S^{-1}(u)|\,du=\int_0^1|R(y)-S(y)|\,dy
\tag{2}
\]
and dominated convergence imply \(\lVert\widehat q_a-q_a\rVert_1\to0\), \(\widehat\mu_a\to\mu_a\), and \(\widehat\sigma_a^2\to\sigma_a^2\).  Since \(|xy-x'y'|\le|x-x'|+|y-y'|\) on \([0,1]^4\), both quantile-product endpoint integrals converge.  Therefore \(\widehat c_e\to c_e\), and the polynomial form of def:moment-matrix gives
\[
\rho_m=\max_e\lVert\widehat Q_e-Q(c_e)\rVert_{\mathrm{op}}\to0
\tag{3}
\]
almost surely.

Put \(K=\Omega^{1/2}\Pi^{-1}\), \(A_e=KQ(c_e)K^{\top}\), and \(\widehat A_e=K\widehat Q_eK^{\top}\).  Exposure positivity makes these matrices well defined and
\[
\lVert\widehat A_e-A_e\rVert_{\mathrm{op}}\le\lVert K\rVert_{\mathrm{op}}^2\rho_m.
\tag{4}
\]
For every published basis, lem:rank-stratified-projector-reduction gives the population and empirical identities \(g_B^+(c_e)=n^{-2}\operatorname{tr}(P_BA_e)\) and \(\widehat g_B(c_e)=n^{-2}\operatorname{tr}(P_B\widehat A_e)\), where \(P_B\) is an orthogonal projector of rank at most two.  Hence trace duality gives
\[
|\widehat g_B(c_e)-g_B^+(c_e)|
\le2n^{-2}\lVert\Omega^{1/2}\Pi^{-1}\rVert_{\mathrm{op}}^2\rho_m.
\tag{5}
\]
No inverse-eigenvalue bound occurs.  Similarly,
\[
|\widehat h_e-h(c_e)|\le n^{-2}\lVert\Pi^{-1}\Omega\Pi^{-1}\rVert_*\rho_m.
\tag{6}
\]
Subtracting (5) from (6) and using the one-Lipschitz property of the maximum proves \(\sup_B|\widehat w(B)-w^+(B)|\le d_m\).  Equation (3) gives \(d_m\to0\).

For arbitrary functions on a common nonempty set, \(|\inf f-\inf g|\le\sup|f-g|\); this proves the value bound without assuming either infimum is attained.  By the definition of infimum, for every \(\varepsilon_m>0\) there exists \(\widehat B_m\) with \(\widehat w(\widehat B_m)\le\widehat v_{\mathrm{mm}}+\varepsilon_m\), and
\[
w^+(\widehat B_m)
\le\widehat w(\widehat B_m)+d_m
\le\widehat v_{\mathrm{mm}}+\varepsilon_m+d_m
\le v_{\mathrm{mm}}^++2d_m+\varepsilon_m.
\tag{7}
\]
Together with \(w^+(\widehat B_m)\ge v_{\mathrm{mm}}^+\), this proves achieved-loss convergence when \(\varepsilon_m\to0\).  The uniform value bound does not create lower semicontinuity at rank changes, so it yields no cross-rank argmin conclusion. \(\square\)